% =============================================================================
%  macros.tex -- mathematical notation
%  Change notation here and it propagates through every section.
% =============================================================================

% --- Number sets ------------------------------------------------------------
\newcommand{\R}{\mathbb{R}}
\newcommand{\N}{\mathbb{N}}
\newcommand{\Z}{\mathbb{Z}}
\newcommand{\C}{\mathbb{C}}

% --- Vectors, matrices, operators -------------------------------------------
\newcommand{\vb}[1]{\bm{#1}}            % vector:  \vb{x}
\newcommand{\mat}[1]{\bm{#1}}           % matrix:  \mat{W}
\newcommand{\T}{\mathsf{T}}             % transpose:  A^\T
\DeclareMathOperator{\Tr}{Tr}
\DeclareMathOperator{\diag}{diag}
\DeclareMathOperator{\rank}{rank}
\DeclareMathOperator{\sign}{sign}
\DeclareMathOperator*{\argmax}{arg\,max}
\DeclareMathOperator*{\argmin}{arg\,min}

% --- Delimiters (starred form auto-sizes: \norm*{...}) ----------------------
\DeclarePairedDelimiter{\abs}{\lvert}{\rvert}
\DeclarePairedDelimiter{\norm}{\lVert}{\rVert}
\DeclarePairedDelimiter{\floor}{\lfloor}{\rfloor}
\DeclarePairedDelimiter{\ceil}{\lceil}{\rceil}
\DeclarePairedDelimiterX{\inner}[2]{\langle}{\rangle}{#1,\,#2}   % \inner{x}{y}
\DeclarePairedDelimiter{\set}{\{}{\}}

% --- Probability and statistics ---------------------------------------------
\DeclareMathOperator{\E}{\mathbb{E}}
\DeclareMathOperator{\Var}{Var}
\DeclareMathOperator{\Cov}{Cov}
\newcommand{\Prob}{\mathbb{P}}
\newcommand{\one}{\mathds{1}}           % indicator function
\newcommand{\given}{\,\middle|\,}       % conditioning bar inside \left...\right
\newcommand{\KL}[2]{D_{\mathrm{KL}}\left(#1\,\middle\|\,#2\right)}
\newcommand{\indep}{\perp\!\!\!\perp}
\newcommand{\iid}{\overset{\text{i.i.d.}}{\sim}}

% --- Calculus ---------------------------------------------------------------
\newcommand{\dd}{\mathrm{d}}            % differential:  \dd x
\newcommand{\pd}[2]{\frac{\partial #1}{\partial #2}}
\newcommand{\od}[2]{\frac{\dd #1}{\dd #2}}
\newcommand{\grad}{\nabla}
\newcommand{\bigO}{\mathcal{O}}

% --- Renormalisation group --------------------------------------------------
\newcommand{\RG}{\mathcal{R}}           % RG map
\newcommand{\coarse}[1]{\tilde{#1}}     % coarse-grained variable
\newcommand{\fixed}[1]{#1^{*}}          % fixed point

% --- Miscellaneous ----------------------------------------------------------
\newcommand{\defeq}{\coloneqq}
\newcommand{\eqdef}{\eqqcolon}
\newcommand{\st}{\text{ s.t. }}
\newcommand{\todo}[1]{\textcolor{red}{\textbf{[TODO: #1]}}}
